% Generated by scripts/assemble.py; corrections belong in source/corrections.json.
\examhead{13}{北京科技大学 2026 年数学分析试卷}{18}
% p019-b12; source PDF page 19
\par\Needspace{5\baselineskip}
\blockanchor{p019-b12}
\majorquestion{一}解答题.\par
% p019-b13; source PDF page 19
\blockanchor{p019-b13}
\question{1}求极限  $\lim_{x\to\left(\frac{\pi}{2}\right)^{-}}\left(\cos x\right)^{\frac{\pi}{2}-x}$ .\par
% p019-b14; source PDF page 19
\blockanchor{p019-b14}
\question{2}设 $a = \sum_{i=1}^{n} a_i, a_i > 0$, 求 $S = \sum_{i=1}^{n} a_i^2$ 的最小值.\par
% p019-b15; source PDF page 19
\blockanchor{p019-b15}
\question{3}求积分 $\int_{\frac{1}{2}}^{\frac{3}{4}} \frac{\arcsin \sqrt{x}}{\sqrt{x(1 - x)}} \mathrm{d}x.$\par
% p019-b16; source PDF page 19
\blockanchor{p019-b16}
\question{4}设  $\{x_{n}\}$  满足  $x_{1}=1, x_{n+1}=x_{n}+\frac{1}{2x_{n}}(n=1,2,\cdots)$ ，求极限  $\lim_{n\to\infty}\frac{x_{n}^{2}-n}{\ln n}$ .\par
% p019-b17; source PDF page 19
\par\Needspace{5\baselineskip}
\blockanchor{p019-b17}
\majorquestion{二}设 $f(x)$ 在 $[0, +\infty)$ 上连续可微，$\lim_{x \to +\infty} f'(x) = A$ (A 为常数)，证明：$\lim_{x \to +\infty} \frac{f(x)}{x} = A$.\par
% p019-b18; source PDF page 19
\par\Needspace{5\baselineskip}
\blockanchor{p019-b18}
\majorquestion{三}已知 $f(x, y) = \varphi(|xy|)$，其中 $\varphi(0) = 0$，且 $u = 0$ 附近满足 $|\varphi(u)| \leq u^2$，证明：$f(x, y)$ 在 (0,0) 处可微。
\majorquestion{四}利用有限覆盖定理证明连续函数的根的存在定理。\par
% p019-b19; source PDF page 19
\par\Needspace{5\baselineskip}
\blockanchor{p019-b19}
\majorquestion{五}求级数 $\sum_{n=1}^{\infty} \frac{(-1)^n a^n}{n \ln(n^2 + n)} x^{2n-1}$ 的收敛域, 其中 $a > 1$.\par
% p019-b20; source PDF page 19
\par\Needspace{5\baselineskip}
\blockanchor{p019-b20}
\majorquestion{六}对于级数 $\sum_{n=1}^{\infty} \frac{\ln(1 + n^2 x)}{n} x^{-(n+1)}$. 证明:\par
% p020-b00; source PDF page 20
\blockanchor{p020-b00}
(1) 该级数在 $[1 + \varepsilon_0, +\infty)(\varepsilon_0 > 0)$ 上一致收敛.\par
% p020-b01; source PDF page 20
\blockanchor{p020-b01}
(2) 该级数在 $(1, +\infty)$ 上不一致收敛.\par
% p020-b02; source PDF page 20
\par\Needspace{5\baselineskip}
\blockanchor{p020-b02}
\majorquestion{七}讨论 $I(a) = \int_{0}^{+\infty} \frac{a}{1 + x^2 a^2} \, \mathrm{d}x$ 在 $[a_0, +\infty)$ 和 $(0, +\infty)$ 的一致收敛性, 其中 $a_0 > 0$.\par
% p020-b03; source PDF page 20
\par\Needspace{5\baselineskip}
\blockanchor{p020-b03}
\majorquestion{八}设 $f(x)$ 是可导函数, 对任意围绕原点一周的定向闭曲线 $L$, 曲线积分 $\int_{L} \frac{x \, \mathrm{d}y - y \, \mathrm{d}x}{f(x) + 4y^2}$ 的值恒为常数 $A$, 且 $f(1) = 2$, 求 $A$ 的值与 $f(x)$.\par
% p020-b04; source PDF page 20
\par\Needspace{5\baselineskip}
\blockanchor{p020-b04}
\majorquestion{九}讨论 $\int_0^{+\infty} \frac{x^{p-1}}{1 + x} \, \mathrm{d}x$ 的敛散性.\par
